\chapter{The Basic SIR Model}

\section{Compartments}

We assume infection is \emph{acute}: an individual is infected, then
is removed (recovers or dies) and stays removed. Three compartments:
\begin{itemize}[nosep]
  \item $S$ --- Susceptible: can become infected.
  \item $I$ --- Infected: currently infectious.
  \item $R$ --- Removed: no longer infectious, not reinfected.
\end{itemize}

%% \begin{center}
%% \begin{tikzpicture}[node distance=2.5cm, >=Stealth]
%%   \node[compartment] (S) {$S$};
%%   \node[compartment] (I) [right=of S] {$I$};
%%   \node[compartment] (R) [right=of I] {$R$};
%%   \draw[->] (S) -- (I);
%%   \draw[->] (I) -- (R);
%% \end{tikzpicture}
%% \end{center}

We write $X,Y,Z$ for the \emph{number} of individuals in each class
(this choice of notation will become clear later). The total
population size $N$ is assumed constant.

\section{Recovery: an exponential process}

Assume individual recovery happens at a constant rate $\gamma$, and
for now ignore new infections. By Basic Probability, a constant recovery rate
means each individual's remaining infectious time $T$ is exponential:
$T\sim\text{Exp}(\gamma)$.

Let $Y_0:=Y(0)$. The probability that an individual is still infected
(not yet recovered) at time $t$ is
\[
P(T>t) = \int_t^\infty \gamma e^{-\gamma s}\,ds = e^{-\gamma t}.
\]
Since each of the $Y_0$ initial infecteds recovers independently with this same law,
\[
Y_t = Y_0\,e^{-\gamma t}.
\]

Differentiating,
\[
\dot Y_t = \frac{d}{dt}\left(Y_0 e^{-\gamma t}\right) = -\gamma Y_0 e^{-\gamma t} = -\gamma Y_t.
\]
So with recovery alone, $\dot Y = -\gamma Y$.

\paragraph{What does $\gamma$ mean?}
$\gamma$ is the per-capita recovery rate. From probability, the
infectious period is $T\sim\text{Exp}(\gamma)$, so the mean infectious
period is $1/\gamma$. Example: if infected individuals recover after 5
days on average, $\gamma = 1/5\ \text{day}^{-1}$.

\section{Adding new infections}

New infections add to $Y$; recoveries remove from $Y$:
\[
\dot Y = (\text{rate of new infections}) - \underbrace{\gamma Y}_{\text{recovery rate}}.
\]

Let us assume for now that there are no recoveries. We derive the rate of new
infections from first principles. Let $\kappa$ denote the number of
contacts per time unit for a susceptible individual. So during a short
time interval this individual has $\kappa \frac{Y}{N} \times \delta t$
contacts with infecteds. We let
\[
c = \Pr(\text{transmission}|\text{contact with infected}).
\]

Hence $1-c$ is the complement (no transmission given a contact) and

\[
\Pr(\text{no infection during time $\delta t$}) = (1-c)^{\kappa \frac{Y}{N} \delta t}
\]

and so

\begin{align*}
  \Pr(\text{infection during time $\delta t$})
  &= 1 - (1-c)^{\kappa \frac{Y}{N} \delta t} \\
  &= 1 - e^{-\beta \frac{Y}{N} \delta t} \quad  (\text{let $\beta:= -\kappa \log(1-c)$})\\
  &= 1 - (1-\beta \frac{Y}{N} \delta t) \quad  (\text{since $e^x \approx 1+x$ for small x})\\
  &= \beta \frac{Y}{N} \delta t  \\
\end{align*}

\begin{definition}
The \emph{force of infection} (FOI) $\lambda_t$ is the per-capita
rate at which susceptibles become infected.
\end{definition}

It is clear that $\lambda_t = \beta \frac{Y}{N}$. This specific form
is called a mass-action force of infection. We can see this by
multiplying by the number of susceptibles $X$:
\begin{align*}
  \text{number of new infections during $\delta t$} &= X \Pr(\text{infection during time $\delta t$}) \\
  &= \beta \frac{Y}{N} \delta t \cdot X
\end{align*}
so the rate of new infections is
\[
\text{rate of new infections} = \frac{\beta \frac{Y}{N} \delta
  t}{\delta t} X = \beta \frac{Y}{N} X.
\]


\paragraph{What does $\beta$ mean?} See below.

\paragraph{What assumptions did we make?}
Contacts do not depend on the infection state of
individuals; the number of contacts between individuals does not
depend on population size $N$: doubling both $N$ and $Y$ leaves
$\lambda_t$ unchanged.


\section{Putting everything together}
We dropped recoveries above only to isolate the derivation of the
infection rate; reinstating the recovery term $\gamma Y$ from
Section~2, we have
\begin{align*}
\dot{X} &= -\text{rate of new infections} = -\beta \frac{Y}{N} X, \\
\dot{Y} &= \text{rate of new infections} - \text{rate of recovery} = \beta\frac{Y}{N}X - \gamma Y, \\
\dot{Z} &= \text{rate of recovery} = \gamma Y.
\end{align*}

Dividing each equation by $N$ (constant), and writing $S=X/N$,
$I=Y/N$, $R=Z/N$ for the population proportions gives the standard
formulation of the SIR model:
\begin{align*}
\dot S &= -\beta S I, \\
\dot I &= \beta S I - \gamma I, \\
\dot R &= \gamma I.
\end{align*}

Note $\dot S+\dot I+\dot R=0$: since $S,I,R$ are proportions,
$S+I+R=1$ is constant. Hence one equation is redundant --- given any
two, the third follows. It is common to drop the $\dot R$ equation and
recover $R=1-S-I$ when needed.


\paragraph{What does $\beta$ mean?}
If $S_0\approx 1$ and $I_0 > 0$ then in one time unit we will have
$\beta SI \approx \beta I$ new infections. So \emph{per individual} we
will have $\beta$ new infections. Thus, $\beta$ is the number of new
infections per individual per time unit that an infected individual
can infect.
